% !TeX program = xelatex
% ЧЕРНОВИК рабочей статьи: система и оценка. Разделы и числа зафиксированы здесь;
% проза намеренно скелетная и помечена % TODO. Сборка: xelatex, затем pandoc в Markdown.

\documentclass[12pt,a4paper]{article}

% ---- Язык ----
\usepackage{polyglossia}
\setmainlanguage{russian}
\setotherlanguage{english}

% ---- Шрифты ----
% Если PT Serif не установлен (sudo apt install fonts-pt),
% замени на "DejaVu Serif" или "Liberation Serif".
\setmainfont{PT Serif}
\setsansfont{PT Sans}
\setmonofont{DejaVu Sans Mono}

% ---- Математика ----
\usepackage{amsmath}
\usepackage{amssymb}

% ---- Таблицы ----
\usepackage{booktabs}
\usepackage{array}

% ---- Геометрия ----
\usepackage{geometry}
\geometry{margin=2.5cm}

% ---- Интервалы ----
\usepackage{parskip}
\usepackage{setspace}
\onehalfspacing

% ---- Ссылки ----
\usepackage[hidelinks]{hyperref}

% ---- Заголовок ----
\title{Ankyra: звуковое нейро-символическое рассуждение\\
через поэтапные разрешимые формализмы\\[0.5em]\large (рабочий черновик)}
\author{Владимир Лапшин\thanks{\texttt{merfill@yandex.ru}}}
\date{}

\begin{document}
\maketitle

\begin{abstract}
Мы представляем Ankyra --- домен-независимый нейро-символический движок рассуждений,
в котором большая языковая модель \emph{предлагает} формализацию, а детерминированный
символический движок \emph{решает} её. Движок покрывает семейство разрешимых
формализмов (определённые хорновские клаузы, стратифицированное отрицание,
позитивная логика первого порядка с дизъюнкцией и кванторами, конечнодоменные
ограничения, точная арифметика и дефолты со специфичностью), каждый --- со звуковой
решающей процедурой и объявленным фрагментом. Каждый ответ несёт механически
проверяемую историю вывода, и ни одно утверждение не входит в теорию без дословной
цитаты или явной пометки гипотезы. Мы приводим proof-of-concept оценку на стандартных
коллекциях: на покрываемых срезах движок достигает высокой точности, порождая
\emph{ноль grounded false proofs}, и честно воздерживается (\texttt{out\_of\_fragment})
там, где его закоммиченная логика не может выразить задачу. Мы не заявляем
state-of-the-art: оценка использует закоммиченные подвыборки, а не полные test-сплиты,
и мы явно фиксируем возникающие статистические и протокольные ограничения.
\end{abstract}

% ======================================================================
\section{Введение}

% TODO: развернуть мотивацию (галлюцинации LLM в рассуждениях; проверяемость;
% граница propose/decide). Явное заявление об области ниже сохранить.

Центральное проектное обязательство: \emph{LLM предлагает, движок решает}. Языковая
модель никогда не владеет истиной. Каждое её утверждение детерминированно
классифицируется как \texttt{derivable}, \texttt{cited} (подкреплено дословной цитатой),
\texttt{hypothesis} (знание, отсутствующее в тексте, с пометкой) или \texttt{rejected}.
Ответ никогда не сильнее своих оснований: вывод из цитат --- \texttt{proven}, вывод,
опирающийся на допущения, --- \texttt{proven\_under(H)}.

\paragraph{Вклад.} % TODO: превратить в список.
\begin{itemize}
	\item Архитектура propose/decide с механически проверяемой историей вывода.
	\item Семейство разрешимых формализмов за \emph{контрактом объявленного фрагмента},
	      с честным воздержанием \texttt{out\_of\_fragment} вместо молчаливого понижения.
	\item Методология гейтования \emph{композитных} бенчмарков, разделяющая потолок
	      охвата (формализм) и потолок извлечения (модель) через gold-fed контрфакт.
	\item Эмпирический proof of concept на стандартных коллекциях с измеряемым
	      инвариантом звуковости (ноль grounded false proofs).
\end{itemize}

\paragraph{Область (явно).} Это proof-of-concept отчёт, а не заявка на
state-of-the-art. Мы оцениваем \emph{закоммиченные подвыборки} стандартных коллекций,
описываем бюджет, не позволяющий прогонять полные test-сплиты, и приводим доверительные
интервалы. Полный охват коллекции недостижим по дизайну для бенчмарков со смешанными
конструкциями, таких как FOLIO (раздел~\ref{sec:results}).

% ======================================================================
\section{Обзор системы}

% TODO: развернуть. Пока держать раздел коротким и фактическим.

Ankyra --- оркестратор формализаций и решающих процедур, а не монолитный решатель.
Задача проходит через: (0)~декомпозицию текста на теорию и запрос; (1)~детерминированное
замыкание теории; (2)~ограниченный цикл предложений (LLM предлагает один типизированный
шаг на волну); и (3)~механическое объяснение, собираемое из истории вывода. Правила,
факты и гипотезы образуют базу, которая только растёт; ответ может меняться, и каждое
изменение записывается.

Детерминированный код никогда не интерпретирует естественный язык. Нормализация
идентификаторов и закрытых enum-полей, которые LLM \emph{уже} извлёк, допустима;
эвристики по сырому тексту, завязанные на слова или фразы-триггеры, --- нет.
Единственные операции над сырым текстом --- это структурные проверки-свидетели
(дословная подстрока/span цитаты, покрытие, длина).

% ======================================================================
\section{Формализмы и решающие процедуры}

Каждая стадия дорожной карты --- именованный формализм со звуковой решающей процедурой,
вводимый только после LLM-free синтетического гейта и гейта на внешней коллекции.

\begin{center}
\begin{tabular}{p{0.8cm}p{5.6cm}p{5.4cm}}
\toprule
\textbf{Стадия} & \textbf{Формализм} & \textbf{Процедура} \\
\midrule
L0 & определённые хорновские клаузы, \texttt{is\_a}, отрицание как комплементарные пары (открытый мир) & semi-naive прямой вывод до неподвижной точки \\
L1 & стратифицированное отрицание / NAF, объявленный закрытый мир & стратифицированное замыкание, ограничения несовместности \\
L2 & позитивная FOL: дизъюнкция, $\exists/\forall$, доказательство разбором случаев & ограниченная резолюция (заземление по конечному домену) \\
L3 & конечнодоменные ограничения (CSP/SAT) & поиск по конечным доменам в репозитории \\
L4 & арифметические термы и уравнения & точное рациональное вычисление / линейное исключение \\
D  & дефолты со специфичностью через \texttt{is\_a} & дефолтный слой (NFA + специфичность) \\
\bottomrule
\end{tabular}
\end{center}

Запрос объявляет требуемый фрагмент; несоответствие отчитывается как
\texttt{out\_of\_fragment}, а не угадывается. L3 и L4 --- отдельные движки за той же
оркестрацией.

% ======================================================================
\section{Методология оценки}
\label{sec:method}

\paragraph{Метрики.} Для каждой коллекции мы приводим \emph{точность по классу}
(трёхклассовая или с выбором варианта) с доверительными интервалами Wilson 95\%.
Гейт также требует \emph{инвариант звуковости}: каждый определённый ответ
\texttt{proven}, и \emph{ноль grounded false proofs} (уверенно неверный определённый
ответ).

\paragraph{Два потолка.} Для композитного бенчмарка мы разделяем:
\emph{охват} (доля gold-формул, которые закоммиченная логика может выразить и решить;
остаток --- честный \texttt{out\_of\_fragment}) и \emph{извлечение} (доля покрытых
строк, которые LLM переводит верно). Контрфакт \emph{gold-fed} подаёт движку
аннотированную формализацию вместо извлечения LLM, отделяя точность метода от точности
языка.

\paragraph{Бюджет и воспроизводимость.} Бесплатного доступа к LLM нет; живые прогоны
платные. Поэтому каждая коллекция входит маленькой закоммиченной подвыборкой,
прогоняется один раз и перепрогоняется только при расхождении. Все LLM-free гейты
(синтетические наборы, gold-fed диагностика) воспроизводимы из одного репозитория.

% ======================================================================
\section{Результаты}
\label{sec:results}

% TODO: проза по подразделам. Числа ниже --- измеренные значения.

\subsection{Синтетические LLM-free гейты}

Эти гейты детерминированы и покрывают каждый конструкт каждого формализма (включая
отрицательные контроли); LLM в цикле нет.

\begin{center}
\begin{tabular}{lcc}
\toprule
\textbf{Гейт} & \textbf{Результат} & \textbf{Grounded false proofs} \\
\midrule
Маршрутизация фрагментов (\texttt{evals.routing\_synthetic}) & 19/19 & 0 \\
L1 отрицание (\texttt{evals.l1\_synthetic}) & 40/40 & 0 \\
L2 клаузальная + равенство (\texttt{evals.l2\_synthetic}) & 65/65 & 0 \\
L3 конечнодоменный CSP (\texttt{evals.l3\_synthetic}) & 31/31 & 0 \\
L4 арифметика (\texttt{evals.l4\_synthetic}) & 37/37 & 0 \\
Дефолты (\texttt{evals.defeasible\_synthetic}) & 8/8 & 0 \\
\bottomrule
\end{tabular}
\end{center}

\subsection{Внешние коллекции}

Таблица~\ref{tab:results} сводит живые/закоммиченные результаты. «Grounded» ---
инвариант звуковости (0 в каждой строке).

\begin{table}[h]
\centering
\begin{tabular}{llrp{2.6cm}}
\toprule
\textbf{Стадия} & \textbf{Коллекция / срез} & \textbf{Результат} & \textbf{95\% CI} \\
\midrule
L0 & ProofWriter, Tier D (OWA) & 297/300 (99,0\%) & [97,1; 99,7] \\
L1 & ProntoQA, tier a & 48/48 (100\%) & [92,6; 100] \\
L1 & ProntoQA, tier b & 160/160 (100\%) & [97,7; 100] \\
L2 & ProntoQA-OOD, tier a & 41/42 (97,6\%) & [87,7; 99,6] \\
L1 & FOLIO, срез отрицания (text-fed) & 11/13 (84,6\%) & [57,8; 95,7] \\
L2 & FOLIO L2 tier a (text-fed, покрытые) & 29/44 (65,9\%) & [51,1; 78,1] \\
L3 & AR-LSAT, eval & 23/30 (76,7\%) & [59,1; 88,2] \\
L4 & GSM8K, eval & 38/40 (95,0\%) & [83,5; 98,6] \\
\bottomrule
\end{tabular}
\caption{Результаты на закоммиченных подвыборках. Знаменатели --- малые выборки, а не
полные test-сплиты; см. раздел~\ref{sec:limitations}.}
\label{tab:results}
\end{table}

\paragraph{ProofWriter (L0).} Tier D: 297/300, 0 grounded false proofs, каждый
определённый ответ \texttt{proven}. Полная коллекция (6\,368 теорий) не прогонялась
(бюджет).

\paragraph{ProntoQA (L1) и ProntoQA-OOD (L2).} ProntoQA tier a 48/48 и tier b
160/160, все \texttt{proven}. ProntoQA-OOD tier a 41/42 (97,6\%), 0 grounded false
proofs.

\paragraph{FOLIO (L1/L2, композитный случай).} Gold-FOL FOLIO охватывает несколько
формализмов сразу, поэтому он прогоняется как внутрифрагментные срезы:
\begin{itemize}
	\item срез L1-отрицания (13 строк): text-fed 11/13, gold-fed 12/13;
	\item L2 tier a (45 строк): охват 44/45 (одна повреждённая аннотация); на покрытых
	      строках text-fed 29/44 против gold-fed 42/44.
\end{itemize}
Разрыв между gold-fed (метод) и text-fed (извлечение) локализует оставшуюся стоимость
в извлечении; потолок охвата на этом срезе фактически закрыт. Дополнение двух
закоммиченных гейт-фрагментов (13 + 93 строки из 204-строчного validation-сплита) несёт
конструкции вне любой закоммиченной стадии и отчитывается как \texttt{out\_of\_fragment}:
xor, $\leftrightarrow$, мультивары, равенство, функциональные термы, схемы аксиом.

\paragraph{AR-LSAT (L3).} Gold-fed 27/27 на реальных играх; живой eval 23/30 (76,7\%),
0 \texttt{grounded\_mismatch}.

\paragraph{GSM8K (L4).} Точная арифметика делает арифметические ошибки невозможными по
построению. Синтетика 37/37; вручную закодированный gold 8/8; живой eval 38/40 (две
неверные строки --- аннотированные ошибки эталона: одна ошибка датасета, одна
неоднозначность).

% ======================================================================
\section{Ограничения}
\label{sec:limitations}

% TODO: это раздел, который обязан оставаться предельно честным. Сохранить целиком.

\paragraph{Размер выборки и статистика.} Каждое число на реальной коллекции опирается
на малую выборку (N = 13, 30, 40, 42, 44, 300), поэтому доверительные интервалы широки
(таблица~\ref{tab:results}); на малых срезах они перекрываются с правдоподобными
baseline'ами и не поддерживают заявку на state-of-the-art. Чтобы поймать разницу около
пяти процентных пунктов при $p\approx0{,}95$ с мощностью $0{,}8$, нужно порядка
$400$ примеров на arm.

\paragraph{Run-once и недетерминизм.} Живой путь вызывает LLM, чей провайдер
недетерминирован даже при нулевой температуре. Прогоны выполняются один раз (политика
бюджета) и повторяются лишь при расхождении, поэтому мы не приводим оценку дисперсии и
не делаем тест на значимость.

\paragraph{Подвыборки, а не test-сплиты.} Мы оцениваем закоммиченные выборки:
ProntoQA-OOD tier a (42 из 5\,800), AR-LSAT eval (30 из $\sim$230), GSM8K eval (40 из
$\sim$1\,319), ProofWriter Tier D (300 из 6\,368 теорий). Полные коллекции не
прогонялись, потому что бесплатного доступа к LLM нет, а прогоны оплачиваются из своего
кармана; это осознанная граница области, а не скрытый провал.

\paragraph{Звуковость относительна формализации.} Движок звуков \emph{относительно}
теории, которую ему подали; перевод естественного языка в эту теорию делает LLM.
Поэтому ошибка извлечения может дать \emph{воздержание} или, в принципе,
\texttt{proven}-ответ о неверной формализации. Мы измеряем \emph{grounded false
proofs} (ноль на каждом отчитанном срезе) и снижаем риски извлечения цитатным
grounding'ом, guard'ом условных цитат, guard'ом доказанного от противного и guard'ами
объявления домена. Это не доказательство истинности end-to-end относительно мира.

\paragraph{У композитного бенчмарка нет единого headline.} Для FOLIO полный охват
коллекции недостижим по дизайну: равенство, функциональные термы, xor, бикондиционал,
мультивары и схемы аксиом лежат вне любой закоммиченной стадии, а полная выводимость
FOL лишь полуразрешима. Поэтому мы приводим стратифицированные срезы и явный счётчик
\texttt{out\_of\_fragment}, а не оценку по всей коллекции.

\paragraph{Нет head-to-head baseline.} Опубликованные числа baseline'ов используют
другие коллекции, сплиты, протоколы и фронтенд-модели, поэтому они не сравнимы
напрямую; мы не запускаем сравнение на общем фронтенде. Вклад --- аргумент о
звуковости и проверяемости, а точность --- поддерживающее свидетельство.

\paragraph{Дрейф модели и провайдера.} Числа привязаны к провайдеру извлечения на
момент прогона; другой или обновлённый провайдер изменит колонки text-fed. Колонки
gold-fed и синтетические --- модельно-независимы.

\paragraph{Дефолты.} Отчитан только синтетический гейт; реальная коллекция
defeasible/NLI остаётся будущей работой.

% ======================================================================
\section{Связанные работы}

% TODO: развернуть с корректным сравнением. Заготовки:
% - Chain-of-thought / self-consistency (рассуждение LLM не сертифицировано).
% - Logic-LM, LINC и другие LLM+symbolic пайплайны (транслятор + решатель).
% - ProofWriter, ProntoQA/ProntoQA-OOD, FOLIO, AR-LSAT, GSM8K как коллекции.
% - Немонотонные / дефолтные рассуждения; история вывода и аудируемый вывод.

% ======================================================================
\section{Заключение}

% TODO: кратко повторить тезис, заявку о звуковости и проверяемости и честную
% область (proof of concept, не state of the art).

Мы представили нейро-символический движок, в котором LLM предлагает, а детерминированный
движок решает, покрывающий семейство разрешимых формализмов за контрактом объявленного
фрагмента. На закоммиченных срезах стандартных коллекций движок достигает высокой
точности, порождая ноль grounded false proofs, и честно воздерживается там, где его
логика не может выразить задачу. Оценка --- proof of concept; ограничения выше
определяют границу того, что поддерживают эти результаты.

\begin{thebibliography}{99}

	\bibitem{ankyra} Ankyra: A domain-agnostic neuro-symbolic reasoning engine. GitHub repository. URL: \url{https://github.com/merfill/ankyra}

	\bibitem{proofwriter} Tafjord O., Mishra B. D., Clark P. ProofWriter: Generating Implications, Proofs, and Abductive Statements over Natural Language. AI2, arXiv:2012.13048.

	\bibitem{prontoqa} Saparov A., He H. Language Models Are Greedy Reasoners: A Systematic Formal Analysis of Chain-of-Thought. ICLR 2023. arXiv:2210.01240.

	\bibitem{prontoqaood} Saparov A., He H. Testing the General Deductive Reasoning Capacity of Large Language Models Using OOD Examples. NeurIPS 2023. arXiv:2305.15269.

	\bibitem{folio} Han S. et al. FOLIO: Natural Language Reasoning with First-Order Logic. arXiv:2209.00840, v3 2024.

	\bibitem{arlsat} Zhong W. et al. AR-LSAT: Investigating Analytical Reasoning of Text. 2021. arXiv:2104.06598.

	\bibitem{agieval} Zhong W. et al. AGIEval: A Human-Centric Benchmark for Evaluating Foundation Models. 2024. arXiv:2304.06364.

	\bibitem{gsm8k} Cobbe K. et al. Training Verifiers to Solve Math Word Problems. 2021. arXiv:2110.14168.

	\bibitem{logiclm} Pan L. et al. Logic-LM: Empowering Large Language Models with Symbolic Solvers for Faithful Logical Reasoning. Findings of EMNLP. 2023.

	\bibitem{linc} Olausson T. X. et al. LINC: A Neurosymbolic Approach for Logical Reasoning by Combining Language Models with First-Order Logic Provers. EMNLP. 2023.

\end{thebibliography}

\end{document}
